\chapter{一张一百块的环球旅行}
\chaptermeta{00}

\begin{ledetext}
先不讲定义。跟着一笔钱走一圈。走完你会发现，你以为自己认识它，其实你只认识它的正面。
\end{ledetext}

\section{周五晚上，98 块}

楼下那家"老陈砂锅"，两个菜一碗汤，扫码，98。

手机叮了一声。你以为故事结束了。

真实发生的事是这样的：某个机房的硬盘上，你名下那行数字从 6432.10 变成 6334.10；老陈名下那行从 21847.60 变成 21945.60。没有一张纸移动过，没有一克金属离开过任何地方，没有人在地下金库里把哪块金砖推到另一格。

两行数字，一减一加。

这就是"付钱"的全部物理过程。

\begin{egbox}{先算一笔账}
如果你觉得"钱当然是纸做的"，那看看这两个数：中国的广义货币 M2 是三百多万亿元，而实际在流通中的现金 M0 只有十几万亿。

剩下那 95\% 以上，\textbf{从生下来就没当过纸}。它们出生是数据库里的一行记录，活着是一行记录，死了是这行记录被划掉。

\end{egbox}

这是要拆的第一层壳。第二层是：这 98 块，接下来去了哪儿。

\section{两周，四个国家}

给它装个 GPS。

\begin{flowlist}
  \flowitem{01}{周六早晨 · 老陈的账}{老陈这个月流水 19 万出头。周一他得给冻品供应商打 3 万 2。钱从他开户的城商行跑到供应商开户的股份行，这两家银行之间怎么结清？在央行的账本上。你吃一顿饭，惊动了中央银行。}
  \flowitem{02}{周一 · 供应商}{供应商收到 3 万 2，转手付掉给养殖场的 2 万 4，剩下的还了经营贷的一期。}
  \flowitem{03}{周一下午 · 一笔钱死了}{还给银行的那部分本金，从这个世界上消失了。不是转手，是消失：银行账上的贷款少一块，存款同步少一块，货币总量真的变少了。这句话现在听着像玄学，第 7 章我会把它拆开摆给你看。}
  \flowitem{04}{周三 · 老陈买了理财}{他账上留 8 万周转，拿 5 万买了一只 R2 的理财。产品说明书第 11 页写着投向：同业存单、政策性金融债、中高等级信用债。其中一小块，最后落到江苏一家做储能电芯的厂子手里，成了它三年期融资的一部分。}
  \flowitem{05}{周四 · 出国}{那家电芯厂用这笔融资订了一台德国的极片分切机。人民币在外汇市场换成欧元。这一刻，你的 98 块参与了一次汇率报价。}
  \flowitem{06}{下周 · 德国的工资单}{德国厂商收到欧元，一部分发工资，一部分进财资账户买了短期债券。全球的钱永远在比价：欧洲短债和美国短债的收益率差上 1 个百分点，就有几十亿资金换个方向。}
  \flowitem{07}{某天早上 · 回到你手机里}{你打开 App，看到基金涨了 0.7\%。}
\end{flowlist}

\begin{pullquote}{}
你花掉的 98 块，两周之内在四个国家转了一圈，顺手参与了一次银行间清算、一次债券配置、一次汇率报价和一次设备投资决策。
\end{pullquote}

这台机器是人类造过的连接最紧密的东西。全天候运转，一处出事全球感冒。而绝大多数每天泡在里面的人，对它怎么转一无所知。

这不怪谁。没人教过。

\section{为什么财经新闻你总是看不懂}

因为新闻只报道\emph{事件}，从不解释\emph{管道}。

\begin{talkbox}
  \talkline{新闻}{"美联储宣布，将联邦基金利率目标区间维持在 3.5\% 至 3.75\% 不变，这是今年连续第五次按兵不动。"}
  \talkline{你}{"哦。"}
  \talkline{翻译}{地球上最重要的那个"借钱价格"没动 → 美国国债收益率的锚没动 → 全世界所有资产在被折算成"今天值多少钱"的时候，用的那个除数没动 → 你的基金净值、你换汇的汇率、你银行理财到期给你的数字、你们公司明年敢不敢扩招，全挂在这根链子上。}
\end{talkbox}

新闻给你一个点。理解世界需要一根线。

这本书只干一件事：把线接上。

\section{不发名词表}

名词表随时能查，也随时会忘。真正稀缺的是因果链：为什么 A 会推出 B，B 又必然带来 C。所以全书从头到尾只讲一条链，反复讲，换着角度讲。

\begin{chainflow}
\chainnode{钱是记账}\chainarrow \chainnode{记账要有人信}\chainarrow \chainnode{信任靠制度和抵押品}\chainarrow \chainnode{银行放贷时创造存款}\chainarrow \chainnode{所以"钱"约等于"债"}\chainarrow \chainnode{债要付息}\chainarrow \chainnode{付息需要增长}\chainarrow \chainnode{增长需要投资}\chainarrow \chainnode{投资需要给风险定价}\chainarrow \chainnode{定价错误加杠杆等于危机}\chainarrow \chainnode{危机后重置资产负债表}\chainarrow \chainnode{循环重来}
\end{chainflow}

你在这条链的末端，是一张很小的资产负债表。你能做的不是预测下一次潮汐什么时候来，而是让这张小表在潮汐里活下来，顺便搭上增长的车。

\section{先扔掉一个东西}

\begin{mythbox}{常见误解}
\mvfrow{误解}{"钱是国家印出来的。政府想多花钱就多印一点，所以才有通货膨胀。"}
\mvfrow{事实}{现代经济里绝大部分货币不是被"印"出来的，是被商业银行"借"出来的。银行批给你 300 万房贷的那一刻，它没从任何人的账户里搬走一分钱，只是在自己账本上同时写下两行：资产端"贷款 300 万"，负债端"你的账户里有 300 万"。\textbf{新的钱就此诞生。}}
这不是什么阴谋论。英格兰银行 2014 年就在官方论文里把它写清楚了，欧洲央行和美联储的表述也一致。是很多教科书没跟上，这是教科书的问题。第 7 章会把它彻底拆开。

\end{mythbox}

\section{四层楼}

越往上越抽象，但每层都踩在下一层的地基上。你可以按顺序爬，也可以直接跳到关心的那层，每章开头会说清它需要什么前置。

\begin{tablewrap}
\begin{xltabular}{\linewidth}{>{\raggedright\arraybackslash}p{0.20\linewidth}XX}
\toprule
\thead{层} & \thead{回答的问题} & \thead{写给谁} \\
\midrule
\textbf{01 · 地基篇} & 钱到底是什么？它凭什么有价值？ & 完全零基础。读完能听懂饭桌上八成的经济话题 \\
\textbf{02 · 管道篇} & 钱怎么被造出来、怎么流动、怎么定价？ & 看得懂新闻标题，但不知道背后的机制 \\
\textbf{03 · 天气篇} & 为什么总有泡沫、危机和周期？ & 想弄明白"为什么每隔十年就来一次" \\
\textbf{04 · 落地篇} & 这些跟我的钱包有什么关系？ & 所有人。尤其是第 23 章 \\
\bottomrule
\end{xltabular}
\end{tablewrap}

\begin{databox}{三条约定}
\textbf{一，遇到英文缩写别慌。}这行黑话密度极高，但每个缩写背后都是一句大白话。本书出现的每个术语都先讲人话，再给名字。

\textbf{二，所有数字都标了时间。}本书写于 2026 年 8 月。金融世界的数字变得很快，结构变得很慢。数字当水位刻度读，结构当河道读。

\textbf{三，不预测。}任何声称能告诉你明年涨跌的人，不是在骗你就是在骗自己。这本书只负责让你看懂机器怎么转，判断留给你。

\end{databox}

\begin{takeaway}{本章一句话}
钱不是纸，是账本上的一行记录。金融不是赌场，是把"今天的钱"和"明天的钱"搬来搬去的一整套管道。

\end{takeaway}

\begin{bridgebox}
\textbf{下一章}既然钱只是一行数字，它凭什么能换来一碗砂锅？这事得从太平洋上一块沉进海里的石头说起。
\end{bridgebox}

